% !TEX root = ../main.tex
\chapter{Abstract}

Quantum computing represents an emerging paradigm with transformative potential for data science, yet a significant gap persists between the theoretical promises of quantum algorithms and their practical applicability on real-world data. This work presents a systematic empirical comparative study of quantum versus classical algorithms across nine use cases covering supervised classification (QSVM), regression (QK-SVR), combinatorial optimisation (QAOA, VQE, fair closed TSP), one-class anomaly detection (OCSVM-ZZ) and one-step lag forecasting. Experiments use three dataset types per case when applicable: map-aligned native data (empirically validated map: \texttt{ZZFeatureMap}, not claimed exclusive), practice data from SKG Tecnología and public benchmarks.

Results show that quantum advantage is \emph{real but conditional}: it appears when the \textbf{feature map} (ZZ or another) is aligned with the data generator; ZZ is the map validated here. QSVM achieves $+38.3$~pp on map-aligned native data; classical methods win on real data without that alignment. In combinatorial optimisation no advantage is observed: QAOA matches the exact optimum in the assignment case (Use Case~3); in the portfolio case (Use Case~4) it reaches the exact optimum in only 6 of 10 seeds (VQE: 3 of 10) whereas the greedy baseline always finds it; and in closed TSP (Use Case~5) QAOA ties the exact optimum with NN+2-opt on Colombia $n=4$ (1867.6~km) under a fair protocol ---modelling value, not speedup. Use Case~7 illustrates the hybrid QAOA--COBYLA paradigm on the minimum dominating set: with $p=3$ the optimum of the 6-node network is recovered by sampling (mean probability of measuring it $5.7\%$, best restart $68.8\%$), tying exact ILP, but only 2 of 5 Erd\H{o}s--R\'enyi instances ($n=6$--$15$) yield a minimum dominating set against 5 of 5 for ILP and Greedy ($-60$~pp), with no measured advantage. Use Cases~8 and~9 confirm the same pattern: OCSVM-ZZ / QK-SVR win only on map-aligned novelty/targets (\texttt{zz\_phase}, \texttt{zz\_lag}); on moons/sine/SKG, classical wins.

Two original contributions are presented. The first is the \emph{quantumness score} $\mathcal{Q}(D)$ (and its hardened revision $\mathcal{Q}_{\mathrm{rob}}$ in QML~v2). The second is the \emph{QCR Framework}, \textbf{formally posed} (protocol + open conjecture): Notebook~09 reported Telco lifts \emph{vs.\ RBF} via re-uploading ($\texttt{reps}=3$); under $\mathcal{Q}_{\mathrm{rob}}$ Telco stays \texttt{classical}, and on SKG rotation only changes the geometry verdict to \texttt{try-ZZ} without QSVM accuracy lift vs.\ GB/RF --- SKG/Telco are \textbf{not} quantum-native (Section~\ref{sec:qml_v2_revision}). Strong validation is post-thesis work. Practical conclusion: adopt quantum algorithms as \emph{selective, complementary tools}. Experiments use Qiskit~2.x (2.3.0 and 2.6.1; see Appendix~\ref{anx:entorno}) on exact statevector simulation.

\vspace{0.5em}
\noindent\textbf{Keywords:} quantum computing, quantum machine learning, QSVM, quantum kernel, QAOA, VQE, combinatorial optimisation, anomaly detection, one-class SVM, variational algorithms, NISQ era, quantum advantage, quantumness score, QCR Framework.

\textbf{Keywords:} quantum computing, quantum advantage, hybrid quantum-classical algorithm, quantumness score, quantum kernel, QAOA, QSVM, OCSVM-ZZ, QCR Framework, minimum dominating set, NISQ
